\bookchapter{第六十二章：介面開發——SwiftUI 入門}

\chapterintro{介面是使用者看到的一切。SwiftUI 是蘋果 2019 年推出的宣告式介面框架——你只需要“宣告”介面長什麼樣，系統自動負責渲染、適配深色模式、適配各種螢幕。到 2026 年，SwiftUI 已是新專案的絕對主流。這一章，咱們學會“畫畫”。}

\booksection{一、SwiftUI 的思想：宣告式，不是命令式}

老式介面開發是“命令式”：先告訴系統“建立按鈕 放在(20,30) 設定顏色”，一步一步指揮。SwiftUI 是“宣告式”：你直接描述“介面上有一個紅色按鈕寫著'點我'”，怎麼畫、什麼時候重畫，系統自己操心。

\begin{lstlisting}[language=Swift]
import SwiftUI

struct ContentView: View {
    var body: some View {      // body：描述"介面長什麼樣"
        VStack {               // 豎著排列
            Text("你好，世界！")
                .font(.largeTitle)      // 修飾符：大標題
            Button("點我") {
                print("被點選了")
            }
            .buttonStyle(.borderedProminent)  // 藍色實心按鈕
        }
        .padding()
    }
}
\end{lstlisting}

關鍵點：檢視 = 遵守 View 協議的結構體，body 屬性描述介面，.修飾符() 鏈式“上妝”。這就是 SwiftUI 的全部基本盤。

\booksection{二、基礎檢視：文字、按鈕、輸入框、圖片}

\begin{center}
\begin{tabularx}{\linewidth}{|X|X|X|}
\hline
\textbf{檢視} & \textbf{幹什麼} & \textbf{示例} \tabularnewline
\hline
Text & 顯示文字 & Text("你好") \tabularnewline
Button & 按鈕 & Button("點我") \{ 動作 \} \tabularnewline
TextField & 輸入框 & TextField("提示", text: \$文字) \tabularnewline
Image & 顯示圖片 & \texttt{Image(systemName: "star.\allowbreak fill")} \tabularnewline
\hline
\end{tabularx}
\end{center}

\begin{center}\small\color{gray}SwiftUI 四大基礎檢視\end{center}

注意 Image(systemName:)——蘋果內建了 SF Symbols 圖示庫（幾千個免費圖示），“star.fill” 這種名字就能直接顯示一個實心星星。這是 iOS 開發的隱藏福利：做介面幾乎不用找圖示素材。

\booksection{三、佈局三件套：VStack、HStack、ZStack}

\begin{itemize}
  \item VStack：豎著排（像排隊打飯）；
  \item HStack：橫著排（像並排坐）；
  \item ZStack：疊著放（像撲克牌摞一起，適合做角標、遮罩）。
\end{itemize}

\begin{lstlisting}[language=Swift]
VStack(alignment: .leading, spacing: 12) {
    HStack {
        Text("帳單")
        Spacer()                 // 彈性空間：把右邊推到盡頭
        Text("-25.5 元").foregroundColor(.red)
    }
    Divider()                    // 分割線
    Text("備註：午飯")
}
.padding(16)                     // 四周留白
.background(Color(.systemGray6)) // 淺灰底
.cornerRadius(12)                // 圓角

// 修飾符順序有意義：先 padding 再 background，背景不含邊距
\end{lstlisting}

\begin{tcbwarn}{注意：修飾符順序有講究}
padding(16).background(灰) 和 background(灰).padding(16) 效果不同——前者“先留白再塗色”（背景不含邊距），後者“先塗色再留白”（背景包含邊距）。新手最常見的“奇怪白邊”都源於此。不確定時調換順序，預覽立刻告訴你答案。

\end{tcbwarn}

\booksection{四、狀態與繫結：讓介面“活”起來}

靜態介面沒意思。SwiftUI 的規則：介面永遠忠實反映狀態，狀態變了，介面自動重畫。

\begin{lstlisting}[language=Swift]
struct CounterView: View {
    @State private var count = 0   // @State：這個值變了要重畫

    var body: some View {
        VStack {
            Text("你點了 \(count) 次")
            Button("點我 +1") {
                count += 1          // 改狀態，介面自動重新整理
            }
        }
    }
}
\end{lstlisting}

輸入框要用“雙向繫結”——\$ 符號把“值和修改權”一起傳進去：

\begin{lstlisting}[language=Swift]
struct InputView: View {
    @State private var text = ""

    var body: some View {
        TextField("輸入點什麼", text: $text)  // $ 雙向繫結
            .textFieldStyle(.roundedBorder)
        Text("你輸入了：\(text)")   // 輸入時實時更新
    }
}
\end{lstlisting}

使用者操作

狀態改變（count += 1）

SwiftUI 感知到變化

自動重畫介面（body 重新計算）

\begin{center}\small\color{gray}圖 3：SwiftUI 的“狀態 介面”單向資料流\end{center}

\booksection{五、預覽、主題與動畫：SwiftUI 的“外掛”}

\begin{itemize}
  \item \#Preview 宏：檔案底部加 \#Preview \{ ContentView() \}，右側畫布實時渲染——改程式碼所見即所得；
  \item 深色模式免費：用系統顏色（Color(.systemGray6)、.primary）而不是寫死顏色，深色模式自動適配，一行不用多寫；
  \item 動畫一行搞定：withAnimation \{ count += 1 \} 或給檢視加 .animation(.default, value: count)——過渡絲滑。
\end{itemize}

\begin{lstlisting}[language=Swift]
struct AnimatedHeartView: View {
    @State private var liked = false

    var body: some View {
        Button {
            withAnimation(.spring(duration: 0.4)) {
                liked.toggle()
            }
        } label: {
            Image(systemName: liked ? "heart.fill" : "heart")
                .font(.system(size: 48))
                .foregroundColor(liked ? .red : .gray)
                .scaleEffect(liked ? 1.3 : 1.0)
        }
    }
}
\end{lstlisting}

\begin{tcbinfo}{小貼士：先功能後動畫}
新手常犯的毛病是“介面還沒做對，就開始糾結動畫”。先讓功能跑通，最後再統一加動效——動畫是錦上添花，功能是雪中送炭。這也是咱們一貫的“簡單至上”原則。

\end{tcbinfo}

\begin{tcbpractice}{來動手：做一個"心情按鈕"}
\booksection{六、常用修飾符速查：一頁在手，介面不愁}

\begin{center}
\begin{tabularx}{\linewidth}{|X|X|}
\hline
\textbf{想做什麼} & \textbf{修飾符} \tabularnewline
\hline
留白 & \texttt{.\allowbreak padding(16)} \tabularnewline
圓角 & \texttt{.\allowbreak cornerRadius(12)} \tabularnewline
背景色 & \texttt{.\allowbreak background(Color(.\allowbreak systemGray6))} \tabularnewline
字型與字號 & \texttt{.\allowbreak font(.\allowbreak title.\allowbreak bold())}/\texttt{.\allowbreak font(.\allowbreak system(size: 20))} \tabularnewline
前景色（文字/圖示） & \texttt{.\allowbreak foregroundColor(.\allowbreak red)} \tabularnewline
撐滿寬度 & \texttt{.\allowbreak frame(maxWidth: .\allowbreak infinity)} \tabularnewline
固定尺寸 & \texttt{.\allowbreak frame(width: 100, height: 44)} \tabularnewline
對齊方式 & \texttt{VStack(alignment: .\allowbreak leading)} \tabularnewline
透明度 & \texttt{.\allowbreak opacity(0.\allowbreak 6)} \tabularnewline
陰影 & \texttt{.\allowbreak shadow(radius: 5)} \tabularnewline
點選手勢 & \texttt{.\allowbreak onTapGesture \{ .\allowbreak .\allowbreak .\allowbreak  \}} \tabularnewline
鍵盤型別 & \texttt{.\allowbreak keyboardType(.\allowbreak decimalPad)} \tabularnewline
\hline
\end{tabularx}
\end{center}

\begin{center}\small\color{gray}SwiftUI 常用修飾符速查\end{center}

記不住沒關係——這頁就是用來查的。用得最多的其實就前四個：padding、圓角、背景、字型。

把模板專案改成一個心情切換器：一個按鈕，點選後在 開心 難過 生氣 之間迴圈切換，文字和背景色跟著變。提示：用陣列存三組（文字+顏色），@State 存下標，點選時 (index + 1) \% 3。做完發給朋友玩，收穫第一波“哇你會做 App 了”。

\end{tcbpractice}

\begin{tcbsummary}{本章小結}
\begin{itemize}
  \item SwiftUI 是宣告式：描述"長什麼樣"，系統負責畫和重畫
  \item 檢視 = 遵守 View 協議的結構體；body 描述介面；修飾符鏈式"上妝"
  \item 四大基礎檢視：Text / Button / TextField / Image；SF Symbols 圖示庫免費管飽
  \item 佈局三件套：VStack 豎排、HStack 橫排、ZStack 疊放；Spacer 是彈性空間
  \item @State 讓介面活起來；\$ 雙向繫結餵給輸入框
  \item \#Preview 所見即所得；系統顏色自動適配深色模式；withAnimation 一行動效
\end{itemize}

\end{tcbsummary}
